\SourceLeaf{399}%
\noindent [oder auch
\[
\begin{aligned}
s\sigma(Y\zeta-Z\eta)&=\alpha(t'P+\theta'Q)-a(t'Q+\theta'R)\\
s\sigma(Z\xi-X\zeta)&=\beta(t'P+\theta'Q)-b(t'Q+\theta'R)\\
s\sigma(X\eta-Y\xi)&=\gamma(t'P+\theta'Q)-c(t'Q+\theta'R).
\end{aligned}
\]
Werden diese Werthe von $Y\zeta-Z\eta,\ Z\xi-X\zeta,\ X\eta-Y\xi$ in dem Ausdrucke für $s\frac{df}{du}$ substituirt, so erhält man vermöge 15:]
\[
\begin{aligned}
ss\sigma\frac{df}{du}={}&-(t''\theta'-\theta''t')(PR-QQ)\\
&+(Pt'+Q\theta')\{(-\tfrac12P''t'+Q't'+\tfrac12R'\theta')t'\\
&\hspace{5em}+(\tfrac12R't'+\tfrac12R''\theta')\theta'\}\\
&-(Qt'+R\theta')\{(\tfrac12P't'+\tfrac12P''\theta')t'\\
&\hspace{5em}+(\tfrac12P''t'+Q''\theta'-\tfrac12R'\theta')\theta'\}.
\end{aligned}
\]
[Mithin ergibt sich schliesslich für $\frac{df}{du}$ die Gleichung]
\begin{equation}\tag*{[16]}
\left\{\begin{aligned}
&(Pt't'+2Qt'\theta'+R\theta'\theta')\BareRoot{(PR-QQ)}.\frac{df}{du}\\
&\quad=(t'\theta''-\theta't'')(PR-QQ)\\
&\qquad+\tfrac12(Pt'+Q\theta')(-P''t't'+2Q't't'+2R't'\theta'+R''\theta'\theta')\\
&\qquad-\tfrac12(Qt'+R\theta')(P't't'+2P''t'\theta'+2Q''\theta'\theta'-R'\theta'\theta'),
\end{aligned}\right.
\end{equation}
[in der ausser $t'=\frac{dt}{du}$, $t''=\frac{ddt}{du^2}$, $\theta'=\frac{d\theta}{du}$, $\theta''=\frac{dd\theta}{du^2}$ nur die Coefficienten $P,Q,R$ von $ss$ und deren erste Ableitungen nach $t$ und $\theta$ vorkommen.]
\Needspace{9\baselineskip}
\begin{center}\bfseries [5.]\end{center}
\noindent\hspace*{1.2em}Um den Ausdruck $\left[\text{für }\frac{df}{du}\right]$ einfacher zu machen, führe man ein [den Winkel $\chi$, den die Tangente der Curve mit der Tangente der Linie $\theta=\text{const.}$ in dem betrachteten Punkte der Fläche bildet. Dann ist:]
\begin{equation}\tag*{[17]}
\left\{\begin{aligned}
\operatorname{tang}\chi&=\frac{\theta'\BareRoot{(PR-QQ)}}{Pt'+Q\theta'}\\
\cos\chi&=\frac{Pt'+Q\theta'}{\BareRoot{P}\BareRoot{(Pt't'+2Qt'\theta'+R\theta'\theta')}}\\
\sin\chi&=\frac{\theta'\BareRoot{(PR-QQ)}}{\BareRoot{P}\BareRoot{(Pt't'+2Qt'\theta'+R\theta'\theta')}}.
\end{aligned}\right.
\end{equation}
[Ferner sei $\omega$ der Winkel zwischen den Linien $t=\text{const.}$ und $\theta=\text{const.}$,
